% Slajd 4 – mapa demonstracji na żywo (pkt 6 wyzwania: określić potrzeby grupy, sprawdzić miejsce,
% pokazać konkretne bariery i udogodnienia, źródło i datę informacji, oznaczenie danych niepełnych
% i przykładowych). Slajd wisi na ekranie tuż przed przejściem do aplikacji; szczegóły w notatkach.
\begin{frame}{Demo: Anna na wózku szuka kawiarni w~centrum}
  \begin{columns}[T,onlytextwidth]
    \begin{column}{0.63\textwidth}
      \scriptsize
      \renewcommand{\arraystretch}{1.1}
      \begin{tabularx}{\textwidth}{@{}>{\bfseries\color{accNavy}}c >{\raggedright\arraybackslash}p{30mm} >{\raggedright\arraybackslash}X@{}}
        \toprule
        & \textbf{Krok} & \textbf{Na co zwrócić uwagę} \\
        \midrule
        1 & Profil potrzeb „Wózek” &
            mapa i~listy pokazują tylko bariery ważne dla wózka; ranking według dopasowania \\
        2 & Pytanie: „kawiarnia w~centrum, na wózku, z~toaletą” &
            ranking z~uzasadnieniem: co miejsce spełnia, a~czego \emph{nie wiadomo} \\
        3 & Karta miejsca &
            każdy fakt ma źródło, datę i~status; brak danych nazwany wprost, z~przyciskiem „Uzupełnij” \\
        4 & Zgłoszenie „Brak windy” na przystanku, waga „blokuje” &
            kwadrat = blokuje; głosy „nadal jest” / „naprawione”; awaria wygasa po 14~dniach \\
        5 & Panel właściciela: „Potwierdź informacje” &
            status „Potwierdzone przez właściciela”; Anna dostaje powiadomienie \\
        \bottomrule
      \end{tabularx}

      \vspace{1.5mm}
      \textbf{Przypadki brzegowe w~demo:} dane sprzeczne $\rightarrow$ „Wymaga ponownej weryfikacji”;
      dane miasta niedostępne $\rightarrow$ snapshot z~datą „stan na”; brak informacji nigdy nie znaczy „dostępne”.

      \vspace{1mm}
      \muted{Dane przykładowe są oznaczone: konta demo (\texttt{anna@}, \texttt{kasia@}, \texttt{admin@accessly.test}),
        zgłoszenia testowe, etykieta „Tryb demo”.}
    \end{column}
    \begin{column}{0.35\textwidth}
      \screenshot{karta-dopasowanie}\par
      \vspace{0.3mm}
      \muted{Krok 3: podsumowanie dopasowania -- co spełnia, czego nie, czego nie wiadomo.}
      \vspace{1.5mm}
      \screenshot{karta-atrybut}\par
      \vspace{0.3mm}
      \muted{Jeden fakt: wartość, źródło (OpenStreetMap), data (2.10.2026), status i~przyciski
        Aktualne / Nieaktualne / Popraw.}
    \end{column}
  \end{columns}

  \note[item]{Czas: ok. 3 minuty. Przed wejściem: aplikacja na telefonie i~laptopie, profil „Wózek” zaznaczony, mapa na Rynku Głównym (zoom 16), panel właścicielki otwarty w~drugiej karcie przeglądarki.}
  \note[item]{Krok 1 (0:00–0:30): kliknąć chip „Wózek”; pokazać toast z~liczbą zgłoszeń i~listę „Najbliżej środka mapy”.}
  \note[item]{Krok 2 (0:30–1:00): zakładka Miejsca, wpisać „Kawiarnia w~centrum, wjadę na wózku, przyda się dostępna toaleta” – w~trybie reguł wynik jest deterministyczny; przeczytać „Spełnia: wejście bez schodów, toaleta; brak danych: szerokie drzwi”.}
  \note[item]{Krok 3 (1:00–1:45): otworzyć kartę; pokazać palcem trzy rzeczy przy atrybucie: źródło, datę, status. Potem „Brak danych” – nasz najważniejszy komunikat: nie zgadujemy. Przyciski Aktualne / Nieaktualne / Popraw, „Uzupełnij” i~pytanie do właściciela.}
  \note[item]{Krok 4 (1:45–2:20): „Zgłoś problem”, pinezka na przystanku, typ „Brak windy”, waga „blokuje”; pokazać kształt znacznika (kwadrat = blokuje) i~głosy. Zgłoszenie można zrobić bez konta.}
  \note[item]{Krok 5 (2:20–3:00): w~drugiej karcie \texttt{/owner.html} jako \texttt{kasia@accessly.test} (Camelot Cafe) kliknąć „Potwierdź informacje”; wrócić do karty – status się zmienił.}
  \note[item]{Przypadki brzegowe (gdy jury pyta lub zostaje czas): historia atrybutu ze sprzecznymi głosami; strona „Źródła danych” z~datami snapshotów; miejsce z~brakującym atrybutem nie dostaje odznaki „Dopasowane do Ciebie”.}
  \note[item]{Jeśli padnie sieć: wszystko poza kafelkami mapy działa z~danych lokalnych. Dane przykładowe: trzy konta demo i~zgłoszenia testowe; widok administratora bez logowania ma etykietę „Tryb demo: dane przykładowe”.}
\end{frame}
